\vspace{0.5em}\noindent\textbf{Mục tiêu Tuần 6:}

\begin{itemize}
\tightlist
\item
  Xây dựng luồng điều phối (Orchestration) tổng hợp tri thức với
  \textbf{LangChain Framework} và mô hình LLM.
\item
  Tích hợp mô hình Ngôn ngữ Lớn (GPT-4o-mini / Bedrock Claude 3.5
  Sonnet) sinh câu trả lời tư vấn streaming real-time.
\item
  Cấu hình hệ thống Giám sát \& Quản trị mô hình (\textbf{LLM
  Observability}) toàn diện với \textbf{Langfuse} và \textbf{AWS
  CloudWatch}.
\end{itemize}

\vspace{0.5em}\noindent\textbf{Các công việc triển khai trong tuần:}

\begingroup
\small
\setlength{\tabcolsep}{3pt}
\renewcommand{\arraystretch}{1.15}
\begin{longtable}[]{@{}
  >{\raggedright\arraybackslash}p{(\columnwidth - 8\tabcolsep) * \real{0.2000}}
  >{\raggedright\arraybackslash}p{(\columnwidth - 8\tabcolsep) * \real{0.2000}}
  >{\raggedright\arraybackslash}p{(\columnwidth - 8\tabcolsep) * \real{0.2000}}
  >{\raggedright\arraybackslash}p{(\columnwidth - 8\tabcolsep) * \real{0.2000}}
  >{\raggedright\arraybackslash}p{(\columnwidth - 8\tabcolsep) * \real{0.2000}}@{}}
\toprule\noalign{}
\begin{minipage}[b]{\linewidth}\raggedright
Thứ
\end{minipage} & \begin{minipage}[b]{\linewidth}\raggedright
Công việc
\end{minipage} & \begin{minipage}[b]{\linewidth}\raggedright
Ngày bắt đầu
\end{minipage} & \begin{minipage}[b]{\linewidth}\raggedright
Ngày hoàn thành
\end{minipage} & \begin{minipage}[b]{\linewidth}\raggedright
Nguồn tài liệu
\end{minipage} \\
\midrule\noalign{}
\endhead
\bottomrule\noalign{}
\endlastfoot
\textbf{2} & - Lập trình luồng điều phối (Orchestration) trong AI
Backend sử dụng \textbf{LangChain Framework}.- Xây dựng QA Prompt chuyên
sâu tổng hợp 4 nguồn thông tin: Câu hỏi tái cấu trúc + Lịch sử chat +
ViMQ Intent/Entities + Top Tài liệu đã Rerank. & 06/07/2026 & 06/07/2026
&
\href{https://python.langchain.com/v0.1/docs/expression_language/primitives/routing/}{RunnableBranch} \\
\textbf{3} & - Tích hợp mô hình LLM (\textbf{GPT-4o-mini / Amazon
Bedrock Claude 3.5 Sonnet}) vào luồng LangChain.- Cấu hình cơ chế truyền
phát trực tiếp (Streaming): Token sinh ra từ LLM đẩy qua gRPC
-\textgreater{} FastAPI Gateway (SSE) về Client. & 07/07/2026 &
07/07/2026 &
\href{https://www.anthropic.com/news/prompt-engineering-for-medical-ai}{Medical
Prompt Safety} \\
\textbf{4} & - Nghiên cứu và triển khai nền tảng \textbf{Langfuse (LLM
Observability)} cho hệ thống Smart Healthcare Platform.- Cấu hình
Langfuse Tracing để ghi nhận toàn bộ vết thực thi của LangChain: thời
gian xử lý từng bước (latency), số lượng token input/output và chi phí
(cost). & 08/07/2026 & 08/07/2026 &
\href{https://www.ncbi.nlm.nih.gov/pmc/articles/PMC7326087/}{Red-Flag
Warnings} \\
\textbf{5} & - Tích hợp hệ thống giám sát hạ tầng với \textbf{AWS
CloudWatch Logs và Metrics}.- Đồng bộ hóa log từ FastAPI Gateway, AI
Backend Microservice và cơ sở dữ liệu NeonDB/RDS lên CloudWatch. &
09/07/2026 & 10/07/2026 &
\href{https://en.wikipedia.org/wiki/Medical_disclaimer}{Medical
Disclaimer} \\
\textbf{6} & - \textbf{Thực hành:} Kiểm thử khả năng quan sát
(Observability Audit) và xử lý ngoại lệ (Error Handling).- Giả lập các
tình huống nghẽn mạng, lỗi gọi API LLM/Reranker và kiểm tra cơ chế
fallback/retry trong LangChain. & 10/07/2026 & 12/07/2026 &
\href{https://platform.openai.com/docs/models/gpt-4o-mini}{GPT-4o-mini
API} \\
\end{longtable}
\endgroup

\vspace{0.5em}\noindent\textbf{Kết quả đạt được:}

\begin{itemize}
\tightlist
\item
  Hoàn thiện hoàn toàn luồng xử lý AI Backend Microservice từ NLP (ViMQ)
  -\textgreater{} RAG (Amazon S3 Vector / Cohere) -\textgreater{} LLM
  Generation (LangChain).
\item
  Đem lại trải nghiệm người dùng xuất sắc với cơ chế streaming token độ
  trễ cực thấp qua gRPC/SSE.
\item
  Đạt tiêu chuẩn quản trị AI chuyên nghiệp với hệ thống Observability
  (Langfuse + AWS CloudWatch), kiểm soát chặt chẽ hiệu năng và chi phí.
\end{itemize}
